\documentclass[11pt,letterpaper]{article}
\usepackage[letterpaper,margin=0.98in,top=0.88in,bottom=0.90in]{geometry}
\usepackage[T1]{fontenc}
\usepackage{lmodern}
\usepackage{microtype}
\usepackage{amsmath,amssymb,amsthm}
\usepackage{fancyhdr}
\usepackage{enumitem}
\usepackage{needspace}

\newtheorem{theorem}{Theorem}
\newtheorem{corollary}{Corollary}
\renewcommand{\qedsymbol}{$\square$}
\setlength{\parindent}{0pt}
\setlength{\parskip}{0.35em}
\setlength{\headheight}{14pt}
\setlength{\headsep}{16pt}
\setlength{\footskip}{24pt}
\setlength{\emergencystretch}{1em}
\pagestyle{fancy}
\fancyhf{}
\renewcommand{\headrulewidth}{0pt}
\fancyfoot[C]{\small\thepage}
\pdfinfo{/Title (Counting Elements of Prime Order: A Note on an Alternative Proof of Cauchy's Theorem) /Author () /Subject () /Keywords ()}

\begin{document}
\begin{center}
{\LARGE\bfseries Counting Elements of Prime Order}\par
\vspace{0.25em}
{\large A Note on an Alternative Proof of Cauchy's Theorem}\par
\end{center}
\vspace{0.48em}

\noindent\textbf{Introduction.}
Cauchy's theorem states that if a prime $p$ divides the order of a finite group $G$, then $G$ contains an element of order $p$. Rotman's \emph{An Introduction to the Theory of Groups} (GTM~148) presents an alternative proof based on cyclically shifting tuples of group elements \cite{Rotman}. Looking more closely at this counting argument, and grouping elements of order $p$ by the cyclic subgroups they generate, yields a formula for their number. We develop this observation and give two applications.

\begin{theorem}\label{thm:main}
Let $G$ be a finite group, and let $p$ be a prime dividing $|G|$. If $N_p$ is the number of elements of $G$ of order exactly $p$, then
\begin{equation}\label{eq:main}
\boxed{N_p=(p-1)(pt+1)\quad\text{for some integer }t\ge 0.}
\end{equation}
Equivalently, the number of subgroups of $G$ of order $p$ is congruent to $1$ modulo $p$.
\end{theorem}

\begin{proof}
Let $1$ denote the identity of $G$, and consider the set
\begin{equation}\label{eq:tuples}
X=\bigl\{(a_1,\ldots,a_p)\in G^p : a_1a_2\cdots a_p=1\bigr\}.
\end{equation}
Any choice of the first $p-1$ entries uniquely determines the last, so
\begin{equation}\label{eq:size}
|X|=|G|^{p-1}\equiv 0\pmod p.
\end{equation}

Cyclic rotation preserves $X$: if $a_1\cdots a_p=1$, then
\[
(a_2\cdots a_p)a_1=a_1^{-1}(a_1\cdots a_p)a_1=1,
\]
and the same reasoning applies to subsequent rotations. The cyclic group of order $p$ therefore acts on $X$ by rotating coordinates. Since $p$ is prime, every orbit has size either $1$ or $p$.

The orbits of size $1$ consist precisely of the constant tuples $(g,\ldots,g)$ with $g^p=1$. As $p$ is prime, the solutions of $g^p=1$ are the identity together with all elements of order $p$. Counting $X$ modulo $p$ thus gives
\begin{equation}\label{eq:congruence}
1+N_p\equiv |X|\equiv0\pmod p,
\qquad\text{so}\qquad N_p\equiv-1\pmod p.
\end{equation}

Every element of order $p$ generates a subgroup of order $p$. Distinct such subgroups intersect only in the identity, so their sets of nonidentity elements partition the set of elements of order $p$ into blocks of size $p-1$. If $r$ is the number of these subgroups, then
\[
N_p=(p-1)r.
\]
Using $p-1\equiv-1\pmod p$ in \eqref{eq:congruence} gives $r\equiv1\pmod p$. Since $r\ge0$, we can write $r=pt+1$ for an integer $t\ge0$. Consequently,
\[
N_p=(p-1)(pt+1),
\]
as claimed. In particular, $r\ge1$, recovering Cauchy's theorem.
\end{proof}

\clearpage
\noindent{\large\bfseries Applications}\par\vspace{0.4em}
\begin{corollary}\label{cor:15}
Every group of order $15$ is cyclic.
\end{corollary}

\begin{proof}
Suppose $|G|=15$ and $G$ is not cyclic. By Lagrange's theorem, every nonidentity element has order $3$ or $5$. Theorem~\ref{thm:main} implies that the number of elements of order $5$ is
\[
N_5=(5-1)(5u+1)=20u+4,\qquad u\ge0.
\]
Since $N_5\le14$, necessarily $u=0$, so $N_5=4$. The other ten nonidentity elements must have order $3$. But Theorem~\ref{thm:main} also requires
\[
N_3=(3-1)(3v+1)=6v+2,\qquad v\ge0,
\]
which cannot equal $10$. This contradiction proves the claim.
\end{proof}

\begin{corollary}\label{cor:pq}
Let $G$ be a group of order $pq$, where $p<q$ are primes. If $p\nmid(q-1)$, then $G$ is cyclic.
\end{corollary}

\begin{proof}
First, Theorem~\ref{thm:main} gives $N_q=(q-1)(qu+1)$ with $u\ge0$. If $u\ge1$, then
\[
N_q\ge(q-1)(q+1)=q^2-1>pq-1,
\]
which is impossible, because $G$ has only $pq-1$ nonidentity elements. Thus $N_q=q-1$.

Suppose $G$ is not cyclic. No element can have order $pq$, so Lagrange's theorem implies that every nonidentity element has order $p$ or $q$. Therefore
\[
N_p=(pq-1)-(q-1)=q(p-1).
\]
Applying Theorem~\ref{thm:main} again, $N_p=(p-1)(pv+1)$ for some $v\ge0$. Dividing by $p-1$ yields $q=pv+1$, hence $p\mid(q-1)$, contrary to the hypothesis. Therefore $G$ is cyclic.
\end{proof}

\begin{thebibliography}{9}
\setlength{\itemsep}{0.2em}
\bibitem{Rotman} J.~J.~Rotman, \emph{An Introduction to the Theory of Groups}, 4th ed., Graduate Texts in Mathematics 148, Springer, 1995, p.~74.
\end{thebibliography}

\end{document}
